% !TeX root = ../../elegantbook-cn.tex

\begin{example}
（2025朝阳期中）已知函数 $f(x)=2\ln x-(2a+1)x+\dfrac{1}{2}ax^2$，其中 $a\in\mathbb{R}$．\\
（1）当 $a=0$ 时，求曲线 $y=f(x)$ 在点 $(1,f(1))$ 处的切线方程；\\
（2）讨论函数 $f(x)$ 的单调性．
\end{example}

\begin{example}
（2025朝阳期中）已知函数 $f(x)=\dfrac{\sin x}{2+\cos x}$．\\
（1）求 $f(x)$ 在区间 $[0,\pi]$ 上的最大值；\\
（2）设 $g(x)=f(x)-\dfrac{1}{\pi}x$，证明：$g(x)$ 在区间 $\left(0,\dfrac{\pi}{2}\right)$ 上有且仅有一个极大值点；\\
（3）当 $x\in(0,\pi)$ 时，直接写出 $\dfrac{x-f(x)}{x+f(x)}$ 的取值范围（结论不要求证明）．
\end{example}

\begin{example}
（2025海淀期中）已知函数 $f(x)=\ln x+\dfrac{a(x-1)}{x-2}$．从下列三个条件中选择一个作为已知条件，使得 $f(x)$ 能够唯一确定，并解决问题：\\
（1）求曲线 $y=f(x)$ 在点 $(1,f(1))$ 处的切线方程；\\
（2）求 $f(x)$ 的单调区间．\\
条件①：$f(1)=0$；条件②：$f(4)=2\ln 2+\dfrac{3}{2}$；条件③：$f'(3)=-\dfrac{2}{3}$．\\
注：如果选择的条件不符合要求，第（1）问得 $0$ 分；如果选择多个符合要求的条件分别解答，按第一个解答计分．
\end{example}

\begin{example}
（2025海淀期中）已知函数 $f(x)=x^3+ax$ 有两个极值点 $x_1,x_2$（$x_1<x_2$），记 $A(x_1,f(x_1))$，$B(x_2,f(x_2))$．\\
（1）若点 $A$ 在直线 $y=-2x$ 上，求 $a$ 的值；\\
（2）若函数 $g(x)=\mathrm{e}^{-ax}$ 的图象上存在点 $P$，使得以 $P$ 为顶点的三角形 $PAB$ 是等腰三角形，求 $a$ 的取值范围．
\end{example}

\begin{example}
（2024海淀期中）已知函数 $f(x)=a\ln(x-a)+\dfrac{1}{2}x^2-(2a+1)x$，其中 $a>0$．\\
（1）若 $f(x)$ 在 $x=4$ 处取得极大值，求 $f(4)$；\\
（2）求函数 $f(x)$ 的零点个数．
\end{example}

\begin{example}
（2024海淀期中）已知函数 $f(x)=\dfrac{x^2-a}{\mathrm{e}^x}$，曲线 $y=f(x)$ 在点 $(0,f(0))$ 处的切线方程为 $y=kx-3$．\\
（1）求 $a,k$ 的值；\\
（2）求函数 $f(x)$ 的最小值．
\end{example}

\begin{example}
（2024朝阳期中）已知函数 $f(x)=ax-\ln(x+1)$，其中 $a\in\mathbb{R}$．\\
（1）当 $a=1$ 时，求 $f(x)$ 的最小值；\\
（2）若 $f(x)$ 存在极小值，求 $a$ 的取值范围．
\end{example}

\begin{example}
（2024朝阳期中）已知函数 $f(x)=\mathrm{e}^x+\cos x$．\\
（1）求曲线 $y=f(x)$ 在点 $(0,f(0))$ 处的切线方程；\\
（2）求函数 $f(x)$ 在区间 $(-\pi,+\infty)$ 上的零点个数；\\
（3）若 $f(m)=n$，且 $m>0$，证明：$n-m>2$．
\end{example}

\begin{example}
（2023朝阳期中）已知函数$f(x) = \mathrm{e}^x -\sin x  -ax^2 (a\in \mathbb{R})$.\\
（1）若$a = 0$，求$f(x)$在区间$[0, \frac{\pi}{2}]$的最小值和最大值.\\
（2）若$a < \frac{1}{2}$，求证：$f(x)$在$x = 0$处取得最小值.
\end{example}

\begin{example}
（2023朝阳期中）已知函数 $f(x)=mx\ln x-x^2+1$，其中 $m\in\mathbb{R}$．\\
（1）当 $m=1$ 时，求曲线 $y=f(x)$ 在点 $(1,f(1))$ 处的切线方程；\\
（2）若 $f(x)\leqslant 0$ 在区间 $[1,+\infty)$ 上恒成立，求 $m$ 的取值范围；\\
（3）比较 $\ln 4$ 与 $\sqrt{2}$ 的大小，并说明理由．
\end{example}

\begin{example}
（2023海淀期中）曲线 $C:y=4-x^2$ 与 $x$ 轴交于不同的两点 $A,B$（点 $A$ 在点 $B$ 的左侧）．点 $P(t,0)$ 在线段 $AB$ 上，且不与端点重合．过点 $P$ 作 $x$ 轴的垂线，交曲线 $C$ 于点 $Q$．\\
（1）若 $\triangle APQ$ 是等腰直角三角形，求其面积；\\
（2）设 $\triangle APQ$ 的面积为 $S(t)$，求 $S(t)$ 的最大值．
\end{example}

\begin{example}
（2023海淀期中）已知函数 $f(x)=\dfrac{\sqrt{x+a}}{x^2+b}$，且 $f(1)=\dfrac14$，$f(4)=\dfrac{2}{19}$．\\
（1）求 $a,b$ 的值；\\
（2）求函数 $f(x)$ 的单调区间；\\
（3）若存在 $k\in\mathbb{R}$，使得直线 $y=kx+m$ 与曲线 $y=f(x)$ 相切，且对任意 $x\in[0,+\infty)$，均有 $kx+m\geqslant f(x)$，求实数 $m$ 的最大值．
\end{example}

\begin{example}
（2022朝阳期中）已知函数 $f(x)=\mathrm{e}^x+a\sin x-1$，其中 $a\in\mathbb{R}$．\\
（1）求曲线 $y=f(x)$ 在点 $(0,f(0))$ 处的切线方程；\\
（2）若 $f(x)$ 在 $x=0$ 处取得极小值，求 $a$ 的值，并说明理由；\\
（3）若存在正数 $m$，使得对任意 $x\in(0,m)$，都有 $f(x)<0$，求 $a$ 的取值范围．
\end{example}

\begin{example}
（2022海淀期中）已知函数 $f(x)=\dfrac13x^3-x^2$．\\
（1）求函数 $f(x)$ 的单调区间；\\
（2）若函数 $f(x)$ 在区间 $(-1,m]$ 上的值域为 $\left[-\dfrac43,0\right]$，求 $m$ 的取值范围．
\end{example}

\begin{example}
（2022海淀期中）已知函数 $f(x)=\mathrm{e}^x-a\sin x$．\\
（1）当 $a=2$ 时，求曲线 $y=f(x)$ 在点 $(0,f(0))$ 处的切线方程；\\
（2）当 $a=1$ 时，证明：函数 $y=f(x)-2$ 在区间 $(0,\pi)$ 上有且仅有一个零点；\\
（3）若对任意 $x\in[0,\pi]$，均有 $f(x)\geqslant 2-\cos x$，求 $a$ 的取值范围．
\end{example}

\begin{example}
（2021朝阳期中）已知函数 $f(x)=\dfrac{ax^2+x-1}{\mathrm{e}^x}$，其中 $a\geqslant 0$．\\
（1）讨论函数 $f(x)$ 的单调性；\\
（2）当 $a>0$ 时，证明：函数 $f(x)$ 在区间 $(0,1)$ 上有且仅有一个零点．
\end{example}

\begin{example}
（2021朝阳期中）已知函数 $f(x)=\tan x-kx^3-x$，其中 $k\in\mathbb{R}$．\\
（1）求曲线 $y=f(x)$ 在点 $(0,f(0))$ 处的切线方程；\\
（2）当 $k=\dfrac13$ 时，证明：对任意 $x\in\left(0,\dfrac\pi2\right)$，都有 $f(x)>0$；\\
（3）若对任意 $x\in\left(0,\dfrac\pi2\right)$，都有 $f(x)>0$，求 $k$ 的最大值．
\end{example}

\begin{example}
（2021海淀期中）已知函数 $f(x)=\ln x$，$g(x)=\dfrac{\mathrm{e}}{x}$．\\
（1）直接写出曲线 $y=f(x)$ 与曲线 $y=g(x)$ 的公共点坐标，并求曲线 $y=f(x)$ 在该点处的切线方程；\\
（2）直线 $x=a$ 分别与曲线 $y=f(x)$、$y=g(x)$ 交于点 $A,B$．当 $a\in(0,\mathrm{e})$ 时，设 $\triangle OAB$ 的面积为 $S(a)$，其中 $O$ 为坐标原点，求 $S(a)$ 的最大值．
\end{example}

\begin{example}
（2021海淀期中）已知函数 $f(x)=x(x^2-3x+a)$，其中 $a\in\mathbb{R}$．\\
（1）当 $a=-9$ 时，求函数 $f(x)$ 的单调递增区间；\\
（2）若函数 $f(x)$ 在区间 $(1,2)$ 上单调递减，求 $a$ 的取值范围；\\
（3）若函数 $f(x)$ 在区间 $(0,2)$ 内有两个极值点 $x_1,x_2$，且 $|f(x_1)-f(x_2)|>|f(x_1)+f(x_2)|$，求 $a$ 的取值范围．
\end{example}
